\section{Giới thiệu và Động lực}
\label{sec:intro}

Tìm kiếm bài báo khoa học là một bài toán khó vì tài liệu liên quan thường phụ thuộc vào thuật ngữ chính
xác, từ viết tắt, thực thể y sinh, và bằng chứng đặc thù cho từng claim. Một hệ thống truy hồi lý tưởng
cần kết hợp cả \emph{exact lexical matching} (từ khóa chính xác) và \emph{semantic matching} (ngữ nghĩa
tương đồng). Trong nhiều năm, pipeline chuẩn để giải quyết bài toán này là: bắt đầu với BM25 (sparse
lexical) và một dense retriever (như SciNCL), kết hợp chúng qua Reciprocal Rank Fusion (RRF), rồi áp dụng
cross-encoder reranking ở tầng cuối để tinh chỉnh thứ hạng~\citep{thakur2021beir,cormack2009rrf}.

Tuy nhiên, pipeline này có ba vấn đề. \emph{Thứ nhất}, nó đắt: BM25, dense retrieval, RRF fusion, và
cross-encoder reranking đều tiêu tốn tài nguyên tính toán. \emph{Thứ hai}, các tầng bổ sung không phải
lúc nào cũng cải thiện chất lượng --- đặc biệt khi dense retriever gốc đã đủ mạnh. \emph{Thứ ba}, khi
fine-tune dense retriever trên một miền chuyên biệt, có sự đánh đổi giữa specialist performance và
generalist transfer: mô hình được fine-tune đạt kết quả cao trên miền đích nhưng suy giảm đáng kể trên
các miền khác. Điều này dẫn đến ba câu hỏi nghiên cứu:

\begin{quote}
\itshape
(1) Liệu một dense retriever hiện đại có thể thay thế toàn bộ pipeline đa tầng cũ với chất lượng tốt hơn
và chi phí thấp hơn? (2) Làm thế nào để tận dụng một specialist fine-tuned trên miền chuyên biệt mà
không làm mất khả năng transfer? (3) Có thể cải thiện transfer bằng các thao tác post-retrieval rẻ chi
phí, không cần GPU inference bổ sung?
\end{quote}

Dự án này trả lời ba câu hỏi trên thông qua bốn đóng góp chính, được xác nhận trên sáu bộ dữ liệu BEIR:

\begin{enumerate}[leftmargin=2em,itemsep=2pt]
  \item \textbf{Chẩn đoán pipeline cũ.} Chúng tôi chỉ ra rằng toàn bộ kiến trúc đa tầng trước đây thực
  chất đang \emph{bù đắp} cho một dense retriever rất yếu: SciNCL ($\ndcg = 0.5640$). Khi thay bằng
  BGE-base-en-v1.5 (109M), kết quả dense đơn thuần đạt $\ndcg = 0.7376$ trên SciFact --- vượt qua
  toàn bộ pipeline cũ (Hybrid RRF + cross-encoder = $0.6939$) với chi phí thấp hơn và độ trễ thấp hơn
  99\%. Cross-encoder làm \emph{giảm} chất lượng có ý nghĩa ($-0.0374$, $p=0.015$) khi áp dụng lên BGE.

  \item \textbf{Fine-tune BGE-small thành SciFact specialist.} Sử dụng recipe self-supervised từ
  vstash~\citep{steffens2026vstash}, chúng tôi fine-tune BGE-small (33M) trên SciFact với RRF
  disagreement signal. Mô hình đạt được $\ndcg = 0.7909$ trên SciFact --- vượt BGE-base (109M, $+0.0533$)
  và vượt vstash reference cùng kích thước ($+0.0202$). Tuy nhiên, transfer average giảm từ $0.3251$
  (BGE-base) xuống $0.3011$ --- một minh họa rõ ràng của specialization-generalization trade-off.

  \item \textbf{Frozen B5 Batch Mode-Switch SGAF.} Chúng tôi đề xuất một controller rẻ chi phí, hoạt
  động ở mức batch, chỉ dùng năm đặc trưng đã có sẵn từ hai lượt retrieval BGE-small/BGE-base (không
  cần inference thêm). Controller tính batch shift score $S$ từ các z-score được chuẩn hóa trên SciFact
  trainfit, và quyết định: nếu $S < 2.0$, giữ specialist-safe mode (BGE-small); nếu $S \ge 2.0$, chuyển
  sang generalist-fallback mode (BGE-base). Frozen B5 đạt transfer average $\ndcg = 0.3249$, gần bằng
  BGE-base ($0.3251$), trong khi vẫn giữ SciFact ở $0.8218$. Toàn bộ quyết định dựa trên thống kê
  frozen từ miền nguồn, không cần label từ target dataset.

  \item \textbf{Frozen P3 Rank-Window Smoothing.} Là một bước hậu xử lý CPU, chỉ kích hoạt trong các
  batch đã được B5 chọn generalist-fallback. P3 pha một prior BGE-small yếu vào top-20 của ranking
  BGE-base qua weighted RRF ($\alpha = 0.10$), cải thiện transfer average lên $0.3293$, vượt BGE-base
  $+0.0042$. Chi phí bổ sung dưới $0.001$\,ms/query. P3 dương trên toàn bộ 6 dataset, significant
  trên NFCorpus ($p=0.02$) và ArguAna ($p=0.02$).
\end{enumerate}

Phần còn lại của báo cáo được tổ chức như sau. Mục~\ref{sec:related} điểm qua các công trình liên quan.
Mục~\ref{sec:method} mô tả kiến trúc pipeline, từ base retrieval đến SGAF. Mục~\ref{sec:setup} trình bày
thiết lập thí nghiệm và chi phí. Mục~\ref{sec:results} báo cáo kết quả chính: base retrieval, SGAF mode
switch, phase evolution, và external validation. Mục~\ref{sec:discussion} thảo luận ý nghĩa. Mục~\ref{sec:limitations}
nêu các hạn chế. Mục~\ref{sec:conclusion} kết luận.
